%-------------------------------------------------------------------------------
\section{(LDAP-specific) Mitigations}
\label{sec:mitigations}


\todo{@Sebastian, delete}


\todo{config of the four authentication mechanisms}
\todo{use TLS}
\todo{prohibit crawling on specific attributes}

%\subsubsection{Valid Reasons for crawlable Servers @Cristoph}
%The question arises if there are reasons for crawlable servers on the public internet.
%- Adressbook of public institutions and postboxes

% RFC Pw Policies DoS: The LDAP directory server SHOULD log as much information as it can (such as client IP address) whenever an account is locked, in order to be able to identify the origin of the attack. Denying anonymous access to the LDAP directory is also a way to restrict this kind of attack. Using the login delay instead of the lockout mechanism will also help avoid this denial of service.

%\subsubsection{Differentiation to webcrawling @Sebastian}

\subsection{Unauthenticated Bind}
\todo{What is this?}

Unauthenticated Bind operations can have significant security issues
   (see Section 6.3.1).  In particular, users intending to perform
   Name/Password Authentication may inadvertently provide an empty
   password and thus cause poorly implemented clients to request
   Unauthenticated access.  Clients SHOULD be implemented to require
   user selection of the Unauthenticated Authentication Mechanism by
   means other than user input of an empty password.  Clients SHOULD
   disallow an empty password input to a Name/Password Authentication
   user interface.  Additionally, Servers SHOULD by default fail
   Unauthenticated Bind requests with a resultCode of
   unwillingToPerform.

SHOULD NOT !!!!! see also security considerations

make statistics about that


Guidelines on securely configuring LDAP servers.
Usage of service accounts for authentication via LDAP
%-------------------------------------------------------------------------------
